\documentclass[10pt,a4paper]{amsart}
\usepackage[margin=19mm]{geometry}
\usepackage{amsmath,amssymb,mathtools}
\usepackage[scr=rsfs,cal=euler]{mathalfa}
\usepackage{xfrac}
\usepackage[unicode,hidelinks,bookmarks=false]{hyperref}
\newcommand{\ie}{i.e.\ }
\newcommand{\cf}{cf.\ }
\newcommand{\resp}{respectively\ }
\newcommand{\Subsection}{\subsection}
\providecommand{\nobreakdash}{\nobreak\hbox{-}}
\setlength{\emergencystretch}{2em}
\multlinegap0pt
\usepackage{tikz}
\usetikzlibrary{cd}
\usepackage{amssymb}
\usepackage{enumerate}
\usepackage{mathtools}
\usepackage{xcolor}
\usepackage[shortlabels]{enumitem}
\newtheorem{lemma}{Lemma}
\renewcommand{\AA}{\mathbb{A}}
\newcommand{\B}{\mathrm{B}}
\newcommand{\C}{\mathbf{C}}
\DeclareMathOperator{\Del}{Del}
\DeclareMathOperator{\Deligne}{Del}
\newcommand{\GG}{\mathbb{G}}
\DeclareMathOperator{\Hod}{Hod}
\newcommand{\Oo}{\mathcal{O}}
\renewcommand{\P}{\mathbf{P}}
\DeclareMathOperator{\Rep}{Rep}
\DeclareMathOperator{\Sim}{Sim}
\newcommand{\TT}{\mathbb{T}}
\DeclareMathOperator{\an}{an}
\newcommand{\ol}[1]{\overline{#1}}
\renewcommand{\to}{\longrightarrow}
\title{Quasi-algebraic quantization for the B-twist Langlands TQFT}
\author{Eric Yen-Yo Chen, Emilio Franco}
\date{Source: arXiv:2609.09098v1 (CC BY 4.0)}
\begin{document}
\maketitle

\noindent\emph{Source and licence.} Quasi-algebraic quantization for the B-twist Langlands TQFT. Version arXiv:2609.09098v1 (CC BY 4.0), distributed by arXiv under the Creative Commons Attribution 4.0 International licence, http://creativecommons.org/licenses/by/4.0/, as recorded in the arXiv licence metadata for this exact version and shown on the abstract page https://arxiv.org/abs/2609.09098v1. Full licence notice: \texttt{licenses/arxiv-2609.09098-CC-BY-4.0.txt}.\par\smallskip

\noindent\textit{Editorial scope.} Counted unit: the article's lemma T\_Del (source0.tex lines 1484--1493). The article's complete original proof of it is delivered with it (source0.tex lines 1495--1531, one \texttt{proof} environment of 37 lines). No mathematics is written, repaired or re-derived: the statement and the proof are the article's own lines, in the article's own notation, and no retained line is substituted or re-flowed.  Retained uncounted as the source-local context the proof needs: lines 1428--1438.  Nothing was omitted from any retained span, so the package carries no omission record.  No retained line contains a citation, so the wrapper carries no bibliography and no bibliography entry.  The retained lines contain the article's own diagram code, which the wrapper typesets with the standard tikz packages; no image file is delivered and the package declares no asset dependency.  Editorial formatting only, in addition: an AMS \texttt{amsart} preamble that restates the article's own declarations and macros used by the retained lines, namely the environments \texttt{lemma} on separate counters, together with the standard packages those lines need.  The retained environments print in this document as Lemma 1 in order of appearance, whereas the article prints the same items with the numbering of its own full document.  The complete native source of the article as distributed in the arXiv e-print for this exact version is bundled unchanged as \texttt{sources/209776/source0.tex}.
\begin{lemma} \label{lm cohomology of the curve}
For $n \in \{ 0, 1, 2 \}$, 
\[
H^n(\C) \cong H^n(C_{\B}) \otimes_k \Oo_{\P^1}(n),
\]
and 
\[
H^n(\C) = 0,
\]
for all other $n$.
\end{lemma}

\begin{lemma} \label{lm description of T_Del}
The tangent complex $\TT_{\Del_{\GG_m}}$ of the Deligne moduli stack for $\GG_m$ has cohomology supported in degrees $-1$, $0$ and $1$, with
\begin{equation} \label{eq description of T_Del}
H^{-1}(\TT_{\Deligne_{\GG_m}(C)}) \cong \tau^*\Oo_{\P^1} \quad \text{and} \quad H^{1}(\TT_{\Deligne_{\GG_m}(C)}) \cong \tau^*\Oo_{\P^1}(2),
\end{equation}
and
\[
H^{0}(\TT_{\Deligne_{\GG_m}(C)}) \cong \tau^*\Oo_{\P^1}(1)^{\oplus 2g} \oplus \tau^*\Oo_{\P^1}(2).
\]
\end{lemma}

\begin{proof}
By construction, $\TT_{\Deligne_{\GG_m}}$ is the homotopy equalizer
\[
\TT_{\Deligne_{\GG_m}} = \mathrm{holim} \left ( \begin{tikzcd}
	\TT_{\Rep^{\an}_{\GG_m} \times \GG_m}   &  \TT_{\Hod_{\GG_m}} \sqcup \TT_{\ol \Hod_{\GG_m}}
	\arrow[shift right, from=1-2, to=1-1]
	\arrow[shift left, from=1-2, to=1-1]
\end{tikzcd} \right ) ,
\]
where the two arrows correspond to $\mu^! \circ (\cdot)^{\an}$ and $\ol\mu^! \circ (\cdot)^{\an}$. Since $\GG_m$ is abelian, the universal $\GG_m$-bundle $pt \to B\GG_m$ induces a trivial adjoint vector bundle. Then, $\TT_{\Sim_{\GG_m}}$ and $\TT_{\Hod_{\GG_m}/\AA^1}$ are provided by the pull-back of the cohomology sheaf of the corresponding Simpson shape of the curve (shifted by $1$). In particular, 
\[
H^n(\TT_{\Hod_{\GG_m}/\AA^1}) \cong \tau_{\Hod}^* R^{n+1}\Gamma(C_{\Hod}), \quad H^n(\TT_{\ol\Hod_{\GG_m}/\AA^1}) \cong \tau_{\ol \Hod}^* R^{n+1}\Gamma(C_{\ol \Hod}),
\]
and
\[
H^n(\TT_{\Rep_{\GG_m}}) \cong \tau_{\B}^* R^{n+1}\Gamma(C_{\B}).
\]
Recalling that $\TT_{\P^1}$ is fully supported in degree $0$, by smoothness, one has for every $n \neq 0$
\[
H^n(\TT_{\Hod_{\GG_m}/\AA^1}) \cong H^n(\TT_{\Hod_{\GG_m}}), \quad H^n(\TT_{\ol\Hod_{\GG_m}/\AA^1}) \cong H^n(\TT_{\ol\Hod_{\GG_m}}), 
\] 
and 
\[
H^n(\TT_{\Rep_{\GG_m} \times \GG_m}) \cong H^n(\TT_{\Rep_{\GG_m}}),
\]
while
\[
H^0(\TT_{\Hod_{\GG_m}}) \cong \tau_{\Hod}^* R^{1}\Gamma(C_{\Hod}) \boxplus \Oo_{\AA^1}, \quad H^0(\TT_{\ol\Hod_{\GG_m}}) \cong \tau_{\ol \Hod}^* R^{1}\Gamma(C_{\ol \Hod}) \boxplus \Oo_{\AA^1},
\]
and
\[
H^n(\TT_{\Rep_{\GG_m}}) \cong \tau_{\B}^* R^{n+1}\Gamma(C_{\B}) \boxplus \Oo_{(\AA^1- 0)},
\]
noting that the extra terms carry a weight $1$ action of $\GG_m$.

Hence, the proof follows immediately from Lemma \ref{lm cohomology of the curve}.
\end{proof}

\end{document}
